\section{Four-Family Combinations and Data-Diversity Scaling}
\label{sec:four-family-diversity}

The preceding experiment evaluated dynamics, rhythm, and phrase/section
features while keeping the high-frequency detector separate.  This section
answers the joint question: all 15 non-empty combinations of four families are
evaluated, and the quantity and generator diversity of the new training data
are varied independently.  The families are denoted by \textbf{S} (corrected
high-frequency Demucs-vocal spectrum), \textbf{D} (non-vocal dynamics),
\textbf{R} (rhythm change), and \textbf{P} (phrase/section bar-length variety).
Code and process references are \coderef{C40}--\coderef{C42} and
\outref{O27}--\outref{O29}; numerical results are \analysisref{A19}--\analysisref{A21}.

\subsection{Change of Evaluation Status and Condition-Disjoint Protocol}

At the user's direction, the 5,500 AIME items are now used to improve and
select the detector.  The former external test is therefore \emph{consumed}:
its earlier frozen evaluation remains a valid historical result for the old
model, but it is no longer an untouched confirmation for the new pipeline.
The present numbers are explicitly exploratory development/CV estimates.

The already frozen five 100-condition groups support five-fold,
condition-disjoint evaluation.  In each fold, one group is held out and the
other 400 prompt conditions form the training pool.  New per-source training
amounts are 25, 50, 100, 200, and 400 conditions.  The same condition IDs are
selected for MTG-Jamendo and every selected AI generator, preserving the
prompt/tag matching.  Across folds every one of the 500 items per source serves
as held-out data once and training data in the other folds.

The 1,600 legacy rows remain in every fit: 400 Human FMA and 400 each from
Suno, HeartMuLa, and ACE-Step.  The number of added AIME AI generator families
is independently set to 1, 2, 4, 6, 8, or 10.  Deterministic cyclic generator
subsets give every generator equal inclusion frequency for each count below
ten.  Binary classes receive equal total weight, and sources receive equal
weight within a class.  All combinations use the same linear ridge score
(penalty 10, decision threshold 0.5), training-median imputation, and explicit
missingness indicators.  No hyperparameter is selected on the reported folds.
The primary metric is generator-macro ROC-AUC: every generator is compared with
the same held-out Human conditions before averaging.

\subsection{High-Frequency Spectral Family}

The spectral family contains 16 scalar features on the 10-second htdemucs
vocal stem: 1--5 and 5--16 kHz tilt; 5--16, 12--20, and 5--10 kHz ratios;
high-frequency flatness, entropy, crest, residual peak density and periodicity;
frame flux and similarity; high-frequency power variation and 4--12 Hz
modulation share; and sibilance contrast and burst rate.  The MUSDB-oracle
average Demucs response is removed by multiplying STFT power by
$10^{-b(f)/10}$.  Frame membership is frozen from the raw-vocal RMS and no
phase-dependent feature is used.  The formal extraction produced 7,100 unique
rows with zero non-finite spectral values in 44.8 seconds.

\subsection{All Combinations by New-Data Quantity}

Table~\ref{tab:four-family-quantity} reports all 15 combinations when all ten
new generator families are represented in training.  The legacy-only column
uses none of the AIME rows.  The remaining columns are five-fold estimates and
therefore stop at 400 training conditions per source.

\begin{table}[H]
\centering
\scriptsize
\setlength{\tabcolsep}{4pt}
\caption{Generator-macro ROC-AUC by feature-family combination and new AIME
training conditions per source.  S=spectral, D=dynamics, R=rhythm, and
P=phrase/section.}
\label{tab:four-family-quantity}
\begin{tabular}{lrrrrrr}
\toprule
Combination & Legacy 0 & 25 & 50 & 100 & 200 & 400 \\
\midrule
D & 0.508 & 0.501 & 0.510 & 0.512 & 0.512 & 0.511 \\
P & 0.500 & 0.500 & 0.500 & 0.500 & 0.500 & 0.500 \\
R & 0.477 & 0.577 & 0.590 & 0.576 & 0.579 & 0.585 \\
S & 0.434 & 0.651 & 0.650 & 0.671 & 0.679 & 0.685 \\
D+P & 0.508 & 0.501 & 0.510 & 0.512 & 0.512 & 0.511 \\
D+R & 0.492 & 0.567 & 0.582 & 0.576 & 0.577 & 0.576 \\
R+P & 0.477 & 0.577 & 0.590 & 0.576 & 0.579 & 0.585 \\
S+D & 0.441 & 0.649 & 0.652 & 0.674 & 0.684 & 0.687 \\
S+P & 0.434 & 0.651 & 0.650 & 0.671 & 0.679 & 0.685 \\
S+R & 0.438 & 0.661 & 0.664 & 0.675 & 0.683 & 0.689 \\
D+R+P & 0.492 & 0.567 & 0.582 & 0.576 & 0.577 & 0.576 \\
S+D+P & 0.441 & 0.649 & 0.652 & 0.674 & 0.684 & 0.687 \\
S+D+R & 0.441 & 0.660 & 0.669 & 0.685 & 0.692 & \textbf{0.695} \\
S+R+P & 0.438 & 0.661 & 0.664 & 0.675 & 0.683 & 0.689 \\
S+D+R+P & 0.441 & 0.660 & 0.669 & 0.685 & 0.692 & \textbf{0.695} \\
\bottomrule
\end{tabular}
\end{table}

The largest change is caused by adapting to new domains and generators, not by
raw sample count.  S+D+R rises from AUC 0.441 under legacy-only training to
0.660 after adding only 25 paired conditions per source.  Increasing 25 to 400
adds a further 0.035.  The CV-leading non-redundant combination is S+D+R at
AUC 0.695 (fold-level SD 0.018) and balanced accuracy 0.645; S alone reaches
0.685, so D and R add only 0.010 in the aggregate.  P has zero complete-feature
coverage at ten seconds.  Every combination containing P consequently matches
the corresponding combination without it; this is an observability failure,
not evidence that musical structure is irrelevant.

\begin{figure}[H]
\centering
\includegraphics[width=\linewidth]{../../four_family_diversity_ablation_20260904/figures/four_family_combination_quantity_table.png}
\caption{All 15 family combinations across training quantities.  Values are
five-fold generator-macro ROC-AUC, except the legacy-only baseline.  This is a
development comparison after the AIME test was consumed, not an external
confirmation.  Data reference: \texttt{A19}.}
\label{fig:four-family-quantity}
\end{figure}

\subsection{Generator Diversity Versus Examples per Generator}

Table~\ref{tab:generator-diversity} holds the per-source amount at 400 and
evaluates S+D+R separately on generator families present and absent from
training.  With one new AI family, the seen-family AUC appears strong at 0.694
while unseen-family AUC is below chance at 0.472.  Unseen transfer rises
monotonically to 0.630 when eight new families are represented.  At eight
training generators, however, increasing the per-source amount from 25 to 400
changes unseen AUC only from 0.618 to 0.630.  Generator diversity therefore has
far greater value than adding more tracks from one or two generators.

\begin{table}[H]
\centering
\caption{S+D+R transfer at 400 new conditions per source.  The ten-generator
case has no unseen AIME generator left to test.}
\label{tab:generator-diversity}
\begin{tabular}{rrrr}
\toprule
New AI generators in training & Seen AUC & Unseen AUC & Unseen BA \\
\midrule
1 & 0.694 & 0.472 & 0.497 \\
2 & 0.686 & 0.516 & 0.525 \\
4 & 0.692 & 0.552 & 0.543 \\
6 & 0.689 & 0.580 & 0.561 \\
8 & 0.689 & 0.630 & 0.601 \\
10 & 0.695 & -- & -- \\
\bottomrule
\end{tabular}
\end{table}

\begin{figure}[H]
\centering
\includegraphics[width=0.93\linewidth]{../../four_family_diversity_ablation_20260904/figures/unseen_generator_diversity_curves.png}
\caption{Unseen-generator AUC versus the number of new AI generator families
represented in training.  Curves compare 25, 100, and 400 paired conditions per
source.  Data reference: \texttt{A20}.}
\label{fig:generator-diversity}
\end{figure}

\subsection{The Aggregate Combination Gain Is Not Uniform}

Table~\ref{tab:four-family-generator} compares the corrected spectral family
with S+D+R after training on all ten new families at 400 conditions/source.
The small macro gain hides large opposing changes: D and R improve AudioLDM 2
Large and the MusicGen variants, but reduce Riffusion and both Stable Audio
variants.  The additional families therefore still behave partly as generator
fingerprints.

\begin{table}[H]
\centering
\small
\caption{Per-generator AUC at 400 new training conditions/source.}
\label{tab:four-family-generator}
\begin{tabular}{lrrr}
\toprule
Generator & S & S+D+R & Change \\
\midrule
AudioLDM 2 Large & 0.701 & 0.785 & $+0.084$ \\
AudioLDM 2 Music & 0.765 & 0.760 & $-0.005$ \\
Riffusion & 0.830 & 0.747 & $-0.082$ \\
MusicGen Medium & 0.690 & 0.733 & $+0.043$ \\
MusicGen Large & 0.693 & 0.724 & $+0.030$ \\
MusicGen Small & 0.650 & 0.722 & $+0.072$ \\
Stable Audio v2 & 0.679 & 0.657 & $-0.022$ \\
Mustango & 0.610 & 0.629 & $+0.018$ \\
Stable Audio v1 & 0.632 & 0.606 & $-0.026$ \\
Udio & 0.601 & 0.588 & $-0.013$ \\
\bottomrule
\end{tabular}
\end{table}

\subsection{Deployment Candidate and Remaining Diversity Gap}

After exploratory selection, an S+D+R candidate was refitted with all 500 AIME
rows per source, together with the 1,600 legacy rows.  This 7,100-row model is
stored in \analysisref{A21}.  Because all AIME rows are consumed, no AIME test
score is attached to that refit.  Its status is a shadow-mode deployment
candidate that requires a new untouched test.

AI-generator diversity has increased substantially, but Human diversity remains
inadequate: the 500 new controls all come from MTG-Jamendo and the legacy Human
source is only FMA.  Moreover, all ten new generator cohorts share AIME
packaging and a common tag-derived prompt system.  AUC 0.695 is therefore an
AIME-domain CV estimate rather than evidence of production-wide AI-music
detection.  The next confirmatory corpus must contain independent Human
catalogues, distinct codec/mastering chains, and generator families absent from
all model-selection stages.
